% GENERATED FILE. Edit canonical lessons, then run npm run generate:books.
% canonical-source-hash: fnv1a64:2564dec29cf50cfb
% canonical-lessons: MW-C263-pujari, MW-C263-mali, MW-C263-teli, MW-C263-dhobi, MW-C263-sunar

\chapter{The Priest, the Gardener and the Goldsmith}
\label{ch:mw-the-priest-the-gardener-and-the-goldsmith}
\hlchaptermodality{263}

\begin{chapteropening}
\textbf{By the end of this chapter:} \emph{I can say what the village's tradesmen do.}

Complete the last lesson of chapter 263: I can say what the village's tradesmen do.
\end{chapteropening}

\section[\texorpdfstring{\mw{पुजारी} (pujārī)}{पुजारी (pujārī)}]{\mw{पुजारी} (pujārī) — a temple priest}

\label{lesson:MW-C263-pujari}



\begin{quote}
Before the new one: say the Marwadi for a village headman, then the Marwadi for a craftsman.
\end{quote}

\subsection*{You'll want to know: \mw{पुजारी}}
\textbf{\mw{पुजारी}} — \emph{pujārī} — \textquotedblleft{}a temple priest\textquotedblright{}.

\begin{quote}
\mw{पुजारी} \mw{मंदिर} \mw{मांय} \mw{पूजा} \mw{करै} \mw{है।}
\end{quote}

\begin{quote}
\emph{pujārī mandir mā̃y pūjā karai hai.}
\end{quote}

\begin{quote}
The priest performs worship in the temple.
\end{quote}

A masculine noun.

\begin{cousinweb}
\textbf{\mw{पुजारी}} (\emph{pujārī}) is inherited from Sanskrit \emph{pūjākara}, \textquotedblleft{}performing worship\textquotedblright{}. CDIAL gives no Marwari form in this entry; it lists Hindi \emph{pujerī}, \textquotedblleft{}temple priest\textquotedblright{} (CDIAL 8317a).
\end{cousinweb}

\begin{grammarlens}[title={What you've built}]
One more word for this chapter.
\end{grammarlens}

\subsection*{Your turn}
\begin{itemize}
\raggedright
  \item \emph{Say it:} \emph{pujārī}
  \item \emph{Say it:} \emph{pujārī}, once more
  \item \emph{Say it:} \emph{pujārī mandir mā̃y pūjā karai hai.} — The priest performs worship in the temple.
  \item \emph{Recall:} say the Marwadi for a village headman, then the Marwadi for a craftsman, then say \emph{pujārī} again
\end{itemize}

\begin{tcolorbox}[breakable,colback=teal!4,colframe=teal!35!black,title={Before you move on}]
What does \mw{पुजारी} mean? (\textquotedblleft{}A temple priest\textquotedblright{}.) Say it once more.
\end{tcolorbox}
\section[\texorpdfstring{\mw{माळी} (māḷī)}{माळी (māḷī)}]{\mw{माळी} (māḷī) — a gardener}

\label{lesson:MW-C263-mali}



\begin{quote}
Before the new one: say the Marwadi for a craftsman, then the Marwadi for a temple priest.
\end{quote}

\subsection*{You'll want to know: \mw{माळी}}
\textbf{\mw{माळी}} — \emph{māḷī} — \textquotedblleft{}a gardener\textquotedblright{}.

\begin{quote}
\mw{माळी} \mw{रूंख} \mw{लगावै} \mw{है।}
\end{quote}

\begin{quote}
\emph{māḷī rū̃kh lagāvai hai.}
\end{quote}

\begin{quote}
The gardener plants trees.
\end{quote}

A masculine noun.

\begin{cousinweb}
\textbf{\mw{माळी}} (\emph{māḷī}) is inherited from Sanskrit \emph{mālin}, \textquotedblleft{}garlanded; a gardener\textquotedblright{}; CDIAL lists Old Marwari \emph{mālī} under it (CDIAL 10094).
\end{cousinweb}

\begin{grammarlens}[title={What you've built}]
One more word for this chapter.
\end{grammarlens}

\subsection*{Your turn}
\begin{itemize}
\raggedright
  \item \emph{Say it:} \emph{māḷī}
  \item \emph{Say it:} \emph{māḷī}, once more
  \item \emph{Say it:} \emph{māḷī rū̃kh lagāvai hai.} — The gardener plants trees.
  \item \emph{Recall:} say the Marwadi for a craftsman, then the Marwadi for a temple priest, then say \emph{māḷī} again
\end{itemize}

\begin{tcolorbox}[breakable,colback=teal!4,colframe=teal!35!black,title={Before you move on}]
What does \mw{माळी} mean? (\textquotedblleft{}A gardener\textquotedblright{}.) Say it once more.
\end{tcolorbox}
\section[\texorpdfstring{\mw{तेली} (telī)}{तेली (telī)}]{\mw{तेली} (telī) — an oil-presser}

\label{lesson:MW-C263-teli}



\begin{quote}
Before the new one: say the Marwadi for a temple priest, then the Marwadi for a gardener.
\end{quote}

\subsection*{You'll want to know: \mw{तेली}}
\textbf{\mw{तेली}} — \emph{telī} — \textquotedblleft{}an oil-presser\textquotedblright{}.

\begin{quote}
\mw{तेली} \mw{तेल} \mw{बेचै} \mw{है।}
\end{quote}

\begin{quote}
\emph{telī tel bechai hai.}
\end{quote}

\begin{quote}
The oil-presser sells oil.
\end{quote}

A masculine noun.

\begin{cousinweb}
\textbf{\mw{तेली}} (\emph{telī}) is inherited from Sanskrit \emph{tailika}, \textquotedblleft{}oilman\textquotedblright{}; CDIAL lists Marwari \emph{telī} under it (CDIAL 5963).
\end{cousinweb}

\begin{grammarlens}[title={What you've built}]
One more word for this chapter.
\end{grammarlens}

\subsection*{Your turn}
\begin{itemize}
\raggedright
  \item \emph{Say it:} \emph{telī}
  \item \emph{Say it:} \emph{telī}, once more
  \item \emph{Say it:} \emph{telī tel bechai hai.} — The oil-presser sells oil.
  \item \emph{Recall:} say the Marwadi for a temple priest, then the Marwadi for a gardener, then say \emph{telī} again
\end{itemize}

\begin{tcolorbox}[breakable,colback=teal!4,colframe=teal!35!black,title={Before you move on}]
What does \mw{तेली} mean? (\textquotedblleft{}An oil-presser\textquotedblright{}.) Say it once more.
\end{tcolorbox}
\section[\texorpdfstring{\mw{धोबी} (dhobī)}{धोबी (dhobī)}]{\mw{धोबी} (dhobī) — a washerman}

\label{lesson:MW-C263-dhobi}



\begin{quote}
Before the new one: say the Marwadi for a gardener, then the Marwadi for an oil-presser.
\end{quote}

\subsection*{You'll want to know: \mw{धोबी}}
\textbf{\mw{धोबी}} — \emph{dhobī} — \textquotedblleft{}a washerman\textquotedblright{}.

\begin{quote}
\mw{धोबी} \mw{नदी} \mw{कनै} \mw{रैवै} \mw{है।}
\end{quote}

\begin{quote}
\emph{dhobī nadī kanai raivai hai.}
\end{quote}

\begin{quote}
The washerman lives near the river.
\end{quote}

A masculine noun.

\begin{cousinweb}
\textbf{\mw{धोबी}} (\emph{dhobī}) is inherited from a reconstructed \emph{*dhauvati}, \textquotedblleft{}washes\textquotedblright{}; CDIAL lists Marwari \emph{dhoṇo}, \textquotedblleft{}washes\textquotedblright{}, and Hindi \emph{dhobī}, \textquotedblleft{}washerman\textquotedblright{}, under it (CDIAL 6886).
\end{cousinweb}

\begin{grammarlens}[title={What you've built}]
One more word for this chapter.
\end{grammarlens}

\subsection*{Your turn}
\begin{itemize}
\raggedright
  \item \emph{Say it:} \emph{dhobī}
  \item \emph{Say it:} \emph{dhobī}, once more
  \item \emph{Say it:} \emph{dhobī nadī kanai raivai hai.} — The washerman lives near the river.
  \item \emph{Recall:} say the Marwadi for a gardener, then the Marwadi for an oil-presser, then say \emph{dhobī} again
\end{itemize}

\begin{tcolorbox}[breakable,colback=teal!4,colframe=teal!35!black,title={Before you move on}]
What does \mw{धोबी} mean? (\textquotedblleft{}A washerman\textquotedblright{}.) Say it once more.
\end{tcolorbox}
\section[\texorpdfstring{\mw{सुनार} (sunār)}{सुनार (sunār)}]{\mw{सुनार} (sunār) — a goldsmith}

\label{lesson:MW-C263-sunar}



\begin{quote}
Before the new one: say the Marwadi for an oil-presser, then the Marwadi for a washerman.
\end{quote}

\subsection*{You'll want to know: \mw{सुनार}}
\textbf{\mw{सुनार}} — \emph{sunār} — \textquotedblleft{}a goldsmith\textquotedblright{}.

\begin{quote}
\mw{सुनार} \mw{चूड़ी} \mw{अर} \mw{नथ} \mw{बेचै} \mw{है।}
\end{quote}

\begin{quote}
\emph{sunār chūṛī ar nath bechai hai.}
\end{quote}

\begin{quote}
The goldsmith sells bangles and nose rings.
\end{quote}

A masculine noun.

\begin{cousinweb}
\textbf{\mw{सुनार}} (\emph{sunār}) is inherited from Sanskrit \emph{suvarṇakāra}, \textquotedblleft{}goldsmith\textquotedblright{}. CDIAL gives no Marwari form in this entry; it lists Hindi \emph{sunār} (CDIAL 13520).
\end{cousinweb}

\begin{grammarlens}[title={What you've built}]
One more word for this chapter.
\end{grammarlens}

\subsection*{Your turn}
\begin{itemize}
\raggedright
  \item \emph{Say it:} \emph{sunār}
  \item \emph{Say it:} \emph{sunār}, once more
  \item \emph{Say it:} \emph{sunār chūṛī ar nath bechai hai.} — The goldsmith sells bangles and nose rings.
  \item \emph{Recall:} say the Marwadi for an oil-presser, then the Marwadi for a washerman, then say \emph{sunār} again
\end{itemize}

\begin{tcolorbox}[breakable,colback=teal!4,colframe=teal!35!black,title={Before you move on}]
What does \mw{सुनार} mean? (\textquotedblleft{}A goldsmith\textquotedblright{}.) Say it once more.
\end{tcolorbox}
